% ==================== 储能效率口径的稳健性检验 ====================
%
% 独立可移植片段。依赖宏包：booktabs（table 2）、tabularx（table 1）、amsmath。
% 不依赖任何自定义宏（如 \kw、\E），不引用外部 \label。
%
% 并入正文时：
%   1. 若作为「模型评价与推广」章下的一节，请把 \section 改为 \subsection，
%      两个 \subsubsection 改为 \subsubsection 或直接去掉。
%   2. 表号、节号会自动并入所在章节，无需手改。
%   3. 末段附注中「问题一的 LP 松弛紧性证明」请改为对应的 \ref 交叉引用。
%
% 集成位置、必须连带修改之处、数据来源见同目录 说明.md。

\section{储能效率口径的稳健性检验}
\label{sec:eff-robust}

题附录 1 仅给出“储能设备的充放电效率为 $90\%$”，未指明该 $90\%$ 是\textbf{单向效率}
还是\textbf{往返效率}。这一区别会改变最优调度，进而改变全部费用结果，
因此有必要明确声明本文的口径，并检验结论对该口径是否敏感。

本文主口径取\textbf{单向效率}，即充电、放电效率各为 $90\%$（往返 $0.81$）；
作为稳健性检验，我们按\textbf{往返效率}口径（往返 $0.90$）重跑了全部四问，
两口径的定义对照见表 \ref{tab:eff-def}。对照口径取对称拆分
$\eta_c=\eta_d=\sqrt{0.9}=0.948683$，使充电侧与放电侧损耗相同且往返效率恰为 $0.90$；
若把损耗全部归于一充或一放一侧，往返效率虽同为 $0.90$，
但充、放电量的统计口径将不对称，故不采用。

\begin{table}[H]
\centering
\caption{储能效率的两种口径定义}
\label{tab:eff-def}
\begin{tabularx}{\textwidth}{@{}l l X c c@{}}
\toprule
口径 & 含义 & 电量递推式 & 单向效率 & 往返效率 \\
\midrule
\textbf{主口径}（本文采用） & 充电、放电效率各 $90\%$
  & $E_t=E_{t-1}+0.9c_t-d_t/0.9$ & 0.900000 & \textbf{0.81} \\
对照口径 & 放电量 $/$ 充电量 $=90\%$
  & $E_t=E_{t-1}+\eta c_t-d_t/\eta$ & 0.948683 & \textbf{0.90} \\
\bottomrule
\end{tabularx}
\end{table}

两口径下四问（问题四含两个子问，共五组结果）的总费用对照见表 \ref{tab:eff-cost}。

\begin{table}[H]
\centering
\caption{两种效率口径下的总费用对照（元）}
\label{tab:eff-cost}
\begin{tabular}{@{}l r r r r@{}}
\toprule
结果集 & {主口径（往返 $0.81$）} & {对照口径（往返 $0.90$）} & 差额 & 变化率 \\
\midrule
问题一     & 35126.95    & 33801.50     & -1325.45    & $-3.77\%$ \\
问题二     & 14389135.08 & 13910066.83  & -479068.25  & $-3.33\%$ \\
问题三     & 14022396.98 & 13558297.17  & -464099.81  & $-3.31\%$ \\
问题四-2   & 15163043.89 & 14684028.55  & -479015.34  & $-3.16\%$ \\
问题四-3   & 15061735.43 & 14579714.65  & -482020.78  & $-3.20\%$ \\
\bottomrule
\end{tabular}
\end{table}

\textbf{差异机制}。往返效率由 $0.81$ 提高到 $0.90$（$+11.1\%$），
“低价充电、高价放电”的往返损失随之减少，套利空间扩大。以问题一为例，
其收益分解中\textbf{光伏余电消纳项两口径完全相同}（$4037.75$ 元），
$1325.45$ 元的差额\textbf{全部来自峰谷价差套利项}（由 $8887.35$ 元增至 $10212.81$ 元）。
这与机理一致：口径只影响电能在时间上搬运的效率，不改变光伏余电的消纳量。

\textbf{差异是真实的模型输出差异}。两口径下最优调度本身不同，
并非对同一调度方案重新计价。问题一 $144$ 个时段中有 $112$ 个时段的充、放电量发生变化，
其中 $18$:00--$18$:10 的购电量由 $531.89$ kWh 降至 $64.53$ kWh；
总体规律是损耗变小后，夜间少充、高价时段多放、中午补电减少，
即储能更依赖跨时段搬移而非提前大量储电。

\textbf{结论的方向不变}。五个结果集的费用降幅高度一致（$3.16\%$--$3.77\%$），
说明这是效率口径本身的影响，而不是某一问的模型不稳定；
主口径下的各项定性结论——多时点调整的价值、$18$ 点更新无额外收益、
电价—净负载联合场景的必要性等——在对照口径下\textbf{均保持不变}。

综上，本文以\textbf{储能侧单向效率各 $90\%$（往返 $0.81$）为主口径}，
正文全部结果与交付文件 \texttt{result1/2/3/4-2/4-3.xlsx} 均对应主口径。
对照口径的完整交付表（题目表 1、表 2、表 3）与逐段明细见支撑材料
\texttt{question*/efficiency/}。\textbf{两种口径下均通过全部逐段核验}
（四问合计 $48\,096$ 段：能量平衡、效率递推、SOC 边界、功率上限、
充放电互斥、跨日连续与初末状态、费用复算）。

\noindent\textbf{附注}：问题一的 LP 松弛紧性证明中构造的等 SOC 扰动，
对\textbf{任意往返效率} $\eta_c\eta_d<1$ 均成立
（扰动量取 $\min\{c_t,\ \eta_c\eta_d d_t,\ g_t/(1-\eta_c\eta_d)\}$），
故该证明不依赖主口径的 $0.81$ 取值，在对照口径下同样有效，仍无需引入二元变量。
